% !TEX program = xelatex
\documentclass[zihao=5,fontset=windows,oneside]{exam-zh}
\usepackage{tikz}
\usetikzlibrary{arrows.meta,calc,patterns}
\input{\WorksheetCommon}
\providecommand{\eu}{\mathrm e}
\providecommand{\ii}{\mathrm i}
\providecommand{\vect}[1]{\symbfit{#1}}
\providecommand{\mat}[1]{\begin{pmatrix}#1\end{pmatrix}}
\providecommand{\circlednumber}[1]{\ifcase#1\or①\or②\or③\or④\fi}
\newcommand{\WorkSpace}{\ifdefined\QuestionCrops\else\vfill\fi}
\newcommand{\WorkPageBreak}{\ifdefined\QuestionCrops\else\clearpage\fi}
\newcommand{\SeriesChoice}{\setchoices{label=(\Alph*)}

\begin{question}[index=1]
曲线 $y=(x+2)\mathrm e^{\frac{1}{x}}$ 的渐近线的条数为
\begin{choices}[columns=4]
\item $0$．
\item $1$．
\item $2$．
\item $3$．
\end{choices}
\end{question}
\QuestionBreak

\begin{question}[index=2]
设函数 $f(x)$ 可导，$F(x)=f(x)[1-|\ln(1+x)|]$，则
\begin{choices}[columns=1]
\item $F(x)$ 在 $x=0$ 处不可导．
\item 当 $F(x)$ 在 $x=0$ 处可导时，$f(0)=0$．
\item 当 $F(x)$ 在 $x=0$ 处可导时，$f(0)=1$．
\item $F'(0)=f'(0)$．
\end{choices}
\end{question}
\QuestionBreak

\begin{question}[index=3]
设 $L$ 为曲线 $\left(x-\dfrac{1}{2}\right)^2+\left(y-\dfrac{1}{2}\right)^2=\dfrac{1}{2}$，
取逆时针方向，
\[
\begin{aligned}
I&=\oint_L 4y\dd x+(x+y)^2\dd y,\\
J&=\oint_L 4x\dd x+(x+y)^2\dd y,\qquad
K=\oint_L 4xy\dd x+(x+y)^2\dd y,
\end{aligned}
\]
则 $I,J,K$ 的大小关系为
\begin{choices}[columns=2]
\item $I<K<J$．
\item $J<K<I$．
\item $I<J<K$．
\item $K<I<J$．
\end{choices}
\end{question}
\QuestionBreak

\begin{question}[index=4]
设 $u=f(x^2,y^2,xyz)$，函数 $z=g(x,y)$ 由方程
$\displaystyle\int_{\mathrm e^{xy}}^z\varphi(\mathrm e^{xy}+z-t)\dd t=xyz$
确定，其中 $f$ 可微，$\varphi$ 连续，且满足 $\varphi(z)-xy\ne0$，
则 $x\dfrac{\partial u}{\partial x}-y\dfrac{\partial u}{\partial y}=$
\begin{choices}[columns=2]
\item $2xf'_1+2yf'_2$．
\item $2xf'_1-2yf'_2$．
\item $2x^2f'_1+2y^2f'_2$．
\item $2x^2f'_1-2y^2f'_2$．
\end{choices}
\end{question}
\QuestionBreak

\begin{question}[index=5]
设 $\symbfit{A}=\begin{pmatrix}1&0&1\\5&4&9\\2&4&6\end{pmatrix}$，
$\symbfit{A}\symbfit{x}=\symbf{0}$ 的全部解构成的向量空间称为 $\symbfit{A}$ 的零空间，\\
$\symbfit{A}\symbfit{x}=\symbf{0}$ 的\mbox{基础解系}所含向量的个数称为 $\symbfit{A}$ 的零空间维数，则
\begin{choices}[columns=1]
\item $\symbfit{A}$ 的零空间为 $k(1,1,-1)^{\mathrm T}$，$k$ 为任意常数．
\item $\symbfit{A}$ 的零空间维数为 $2$．
\item $\symbfit{A}^{\mathrm T}$ 的零空间为 $k(1,1,-1)^{\mathrm T}+l(-1,1,0)^{\mathrm T}$，$k,l$ 为任意常数．
\item $\symbfit{A}^{\mathrm T}$ 的零空间维数为 $2$．
\end{choices}
\end{question}
\QuestionBreak

\begin{question}[index=6]
设 $\symbfit{A}$ 为 $3$ 阶非零矩阵，$\symbfit{A}^*$ 是 $\symbfit{A}$ 的伴随矩阵，
若 $\symbfit{A}^*+2\symbfit{A}=\symbfit{O}$，则 $\symbfit{A}^{10}=$
\begin{choices}[columns=2,label-pos=left]
\item $\begin{pmatrix}-4^5&&\\&-4^5&\\&&-4^5\end{pmatrix}$．
\item $\begin{pmatrix}-4^5&&\\&-4^5&\\&&4^5\end{pmatrix}$．
\item $\begin{pmatrix}-4^5&&\\&4^5&\\&&4^5\end{pmatrix}$．
\item $\begin{pmatrix}4^5&&\\&4^5&\\&&4^5\end{pmatrix}$．
\end{choices}
\end{question}
\QuestionBreak

\begin{question}[index=7]
设 $3$ 阶矩阵 $\symbfit{A},\symbfit{B}$ 满足
$\symbfit{A}^2+\symbfit{B}^2=\symbfit{A}\symbfit{B}+\symbfit{B}\symbfit{A}$，
且 $\symbfit{A}\ne\symbfit{B}$，则下列结论错误的是
\begin{choices}[columns=2]
\item $\symbfit{A}-\symbfit{B}$ 的秩为 $1$．
\item $\symbfit{A}-\symbfit{B}$ 的迹为 $0$．
\item $\symbfit{A}-\symbfit{B}$ 可有非零特征值．
\item $\symbfit{A}-\symbfit{B}$ 不可相似对角化．
\end{choices}
\end{question}
\QuestionBreak

\begin{question}[index=8]
设 $X_1,X_2,\ldots,X_{10}$ 是来自总体 $N(0,\sigma^2)$（$\sigma>0$）的简单随机样本，记
\[
\overline X=\frac{1}{10}\sum_{i=1}^{10}X_i,\qquad
S_1^2=\frac{1}{9}\sum_{i=1}^{10}(X_i-\overline X)^2,\qquad
S_2^2=\frac{1}{10}\sum_{i=1}^{10}X_i^2,\quad\text{则}
\]
\begin{choices}[columns=2,label-pos=left]
\item $\dfrac{\sqrt{10}\,\overline X}{S_1}\dist t(10)$．
\item $\dfrac{\sqrt{10}\,\overline X}{S_2}\dist t(10)$．
\item $\dfrac{9S_1^2}{\sigma^2}\dist\chi^2(9)$．
\item $\dfrac{10S_2^2}{\sigma^2}\dist\chi^2(9)$．
\end{choices}
\end{question}
\QuestionBreak

\begin{question}[index=9]
设 $X$ 与 $Y$ 为具有二阶矩的随机变量，且 $DX\ne0$，
若记 $G(a,b)=E\bigl\{[Y-(a+bX)]^2\bigr\}$，则 $G(a,b)$ 在驻点处的值为
\begin{choices}[columns=2]
\item $(1-\rho_{XY}^2)DY$．
\item $(1+\rho_{XY}^2)DY$．
\item $(1-\rho_{XY}^2)\operatorname{Cov}(X,Y)$．
\item $(1+\rho_{XY}^2)\operatorname{Cov}(X,Y)$．
\end{choices}
\end{question}
\QuestionBreak

\begin{question}[index=10]
设总体 $X$ 的分布为
\[
\begin{gathered}
P\{X=k\}=\theta^{k-1}(1-\theta),\quad k=1,2,\ldots,r,\\
P\{X=r+1\}=\theta^r,\quad \theta\in(0,1).
\end{gathered}
\]
从总体中抽取容量为 $n$ 的样本 $x_1,x_2,\ldots,x_n$，
其中有 $M$ 个取值为 $r+1$，\\
则 $\theta$ 的最大似然估计值为
\begin{choices}[columns=2,label-pos=left]
\item $\hat\theta=\dfrac{\displaystyle\sum_{i=1}^n x_i+n}{\displaystyle\sum_{i=1}^n x_i+M}$．
\item $\hat\theta=\dfrac{\displaystyle\sum_{i=1}^n x_i-n}{\displaystyle\sum_{i=1}^n x_i-M}$．
\item $\hat\theta=\dfrac{\displaystyle\sum_{i=1}^n x_i^2+n}{\displaystyle\sum_{i=1}^n x_i^2+M}$．
\item $\hat\theta=\dfrac{\displaystyle\sum_{i=1}^n x_i^2-n}{\displaystyle\sum_{i=1}^n x_i^2-M}$．
\end{choices}
\end{question}
\QuestionBreak
}
\newcommand{\SeriesFill}{%
  \ifdefined\QuestionCrops\else
    \WorksheetTitle{填空题}
    \section{填空题：第11～16题，每小题5分，共30分。}
    \examsetup{question/bottom-sep=11pt}
  \fi
  % !TEX program = xelatex
% 张宇考研数学预测8套卷（2027年，数学一，卷一），原试卷第4页。
\documentclass[zihao=5,fontset=windows,oneside]{exam-zh}
\examsetup{
  page/size=a4paper,
  page/show-head=false,
  page/foot-content={第;页（共;页）},
  sealline/show=false,
  question={hang=true,show-answer=false,top-sep=5pt,bottom-sep=3pt},
  choices={columns=0,column-sep=1.6em,linesep=5pt,label-pos=top-left},
  fillin={no-answer-type=none,width=19mm},
}
\linespread{1.1}
\setCJKfamilyfont{examcoverfont}[BoldFont=FandolSong-Bold.otf]{FandolSong-Regular.otf}
\providecommand{\dd}{\,\mathrm d}
\newcommand{\WorksheetTitle}[1]{%
  \begin{center}
    {\CJKfamily{examcoverfont}\bfseries\fontsize{17}{22}\selectfont 2027年张宇考研数学预测8套卷}\\[5pt]
    {\CJKfamily{examcoverfont}\fontsize{13}{18}\selectfont 数学一 · 第1套 · #1}
  \end{center}\vspace{8pt}%
}
\newcommand{\QuestionBreak}{}
\ifdefined\QuestionCrops
  \geometry{textwidth=468bp,hcentering}
  \renewcommand{\QuestionBreak}{\clearpage}
  \AtBeginDocument{\pagestyle{empty}}
\fi

\AtBeginDocument{\hypersetup{pdftitle={2027年张宇考研数学预测8套卷 · 数学一 · 第1套 · 填空题}}}
\begin{document}

\ifdefined\QuestionCrops\else
  \WorksheetTitle{填空题}
  \section{填空题：第11～16题，每小题5分，共30分。}
  \examsetup{question/bottom-sep=11pt}
\fi

\begin{question}[index=11]
$\displaystyle\lim_{x\to 0}\left(\frac{\mathrm e^x-1}{\sin x}\right)^{\cot x}=$\fillin。
\end{question}
\QuestionBreak

\begin{question}[index=12]
设连续函数 $y=f(x)$ 在区间 $[0,2]$ 上满足 $f(x)+f(2-x)=0$，\par
则 $\displaystyle\int_0^\pi f(1+\cos x)\dd x=$\fillin。
\end{question}
\QuestionBreak

\begin{question}[index=13]
曲线 $\begin{cases}x=t-\ln(1+t^2),\\ y=\arctan t\end{cases}$ 在 $t=2$ 处的曲率为\fillin。
\end{question}
\QuestionBreak

\begin{question}[index=14]
设上半球壳 $x^2+y^2+z^2=R^2\ (z\geq 0,R>0)$ 的面密度
$\rho(x,y,z)=\sqrt{x^2+y^2+z^2}$，则其对\mbox{原点}的转动惯量为\fillin。
\end{question}
\QuestionBreak

\begin{question}[index=15]
已知3阶矩阵 $\symbfit A$ 的每行元素之和均为0，$r(\symbfit A)=2$，\par
则线性方程组 $\symbfit A\symbfit x=\symbfit 0$ 的通解为\fillin。
\end{question}
\QuestionBreak

\begin{question}[index=16]
设随机变量 $X$ 和 $Y$ 均服从 $B\!\left(1,\frac12\right)$，且 $D(X+Y)=1$，\par
则 $X$ 和 $Y$ 的相关系数 $\rho=$\fillin。
\end{question}

\end{document}
%
}
\newcommand{\SeriesLong}{%
  \examsetup{question={show-points=true,points-prelabel={（本题满分},points-postlabel={分）},points-separate-par=true}}
  \ifdefined\QuestionCrops\else
    \WorksheetTitle{解答题}
    \section{解答题：第17～22题，共70分。解答应写出文字说明、证明过程或演算步骤。}
  \fi
  % !TEX program = xelatex
% 张宇考研数学预测8套卷（2027年，数学一，卷一），原试卷第5～8页。
\documentclass[zihao=5,fontset=windows,oneside]{exam-zh}
\examsetup{
  page/size=a4paper,
  page/show-head=false,
  page/foot-content={第;页（共;页）},
  sealline/show=false,
  question={hang=true,show-answer=false,top-sep=5pt,bottom-sep=3pt},
  choices={columns=0,column-sep=1.6em,linesep=5pt,label-pos=top-left},
  fillin={no-answer-type=none,width=19mm},
}
\linespread{1.1}
\setCJKfamilyfont{examcoverfont}[BoldFont=FandolSong-Bold.otf]{FandolSong-Regular.otf}
\providecommand{\dd}{\,\mathrm d}
\newcommand{\WorksheetTitle}[1]{%
  \begin{center}
    {\CJKfamily{examcoverfont}\bfseries\fontsize{17}{22}\selectfont 2027年张宇考研数学预测8套卷}\\[5pt]
    {\CJKfamily{examcoverfont}\fontsize{13}{18}\selectfont 数学一 · 第1套 · #1}
  \end{center}\vspace{8pt}%
}
\newcommand{\QuestionBreak}{}
\ifdefined\QuestionCrops
  \geometry{textwidth=468bp,hcentering}
  \renewcommand{\QuestionBreak}{\clearpage}
  \AtBeginDocument{\pagestyle{empty}}
\fi

\examsetup{question={show-points=true,points-prelabel={（本题满分},points-postlabel={分）},points-separate-par=true}}
\AtBeginDocument{\hypersetup{pdftitle={2027年张宇考研数学预测8套卷 · 数学一 · 第1套 · 解答题}}}
\newcommand{\WorkSpace}{\ifdefined\QuestionCrops\else\vfill\fi}
\newcommand{\WorkPageBreak}{\ifdefined\QuestionCrops\else\clearpage\fi}
\begin{document}

\ifdefined\QuestionCrops\else
  \WorksheetTitle{解答题}
  \section{解答题：第17～22题，共70分。解答应写出文字说明、证明过程或演算步骤。}
\fi

\begin{question}[index=17,points=10]
求函数 $f(x,y)=(y-x)(y-x^2)$ 的极值。
\end{question}
\QuestionBreak\WorkSpace

\begin{question}[index=18,points=12]
已知函数 $g(x)$ 连续，$g(0)=0$，设 $\displaystyle f(x)=\int_0^x g(x^2t)\dd t$，求 $f'(x)$ 的表达式，并判断 $f'(x)$ 在 $x=0$ 处的连续性。
\end{question}
\QuestionBreak\WorkSpace\WorkPageBreak

\begin{question}[index=19,points=12]
若曲线 $y=f(x)$ 的切线斜率在 $[a,b]$ 上单调递增，证明曲线 $y=f(x)$ 在 $[a,b]$ 上是凹的。
\end{question}
\QuestionBreak\WorkSpace

\begin{question}[index=20,points=12]
若对第一卦限内任一有向闭曲面 $\Sigma$ 恒有
\[
  \oiint_{\Sigma} x^2z\varphi(y)\dd y\dd z
  +2xz\varphi'(y)\dd z\dd x-xz^2\mathrm e^y\dd x\dd y=0,
\]
其中 $\varphi(y)$ 具有二阶连续导数，
$\displaystyle\varphi(\pi)=\varphi'(\pi)=\frac{\mathrm e^\pi-1}{2}$，求 $\varphi(y)$。
\end{question}
\QuestionBreak\WorkSpace\WorkPageBreak

\begin{question}[index=21,points=12]
设 $\displaystyle
  \symbfit\alpha_1=\begin{pmatrix}2\\2\\-1\end{pmatrix},\quad
  \symbfit\alpha_2=\begin{pmatrix}-1\\-3\\2\end{pmatrix},\quad
  \symbfit\alpha_3=\begin{pmatrix}0\\-4\\3\end{pmatrix}$，并记
$\displaystyle\symbfit A=\begin{pmatrix}2&-1&0\\2&-3&-4\\-1&2&3\end{pmatrix}$。
\begin{enumerate}[label=（\arabic*）,leftmargin=2.5em,labelsep=.4em,itemsep=3pt,topsep=3pt]
\item 求 $\symbfit\alpha_1,\symbfit\alpha_2,\symbfit\alpha_3$ 的一个极大线性无关组；
\item 求 $\symbfit A^n\ (n>2)$。
\end{enumerate}
\end{question}
\QuestionBreak\WorkSpace

\begin{question}[index=22,points=12]
设随机变量 $X$ 的概率分布为 $\displaystyle P\{X=k\}=\frac{1}{2^{k+1}}+\frac{a}{3^k}$，$k=1,2,\dots$。
\begin{enumerate}[label=（\arabic*）,leftmargin=2.5em,labelsep=.4em,itemsep=3pt,topsep=3pt]
\item 求 $a$ 的值；
\item 求 $X$ 的数学期望 $EX$ 与方差 $DX$；
\item 比较 $E(X\mid X>1)$ 与 $E(X\mid X>2)$ 的大小，并说明理由。
\end{enumerate}
\end{question}
\WorkSpace

\end{document}
%
}
\begin{document}
\csname Series\WorksheetKind\endcsname
\end{document}
